\appendix
\section{Supplementary protocol details}

\subsection{Reproducible audit sequence}

For each primary-parameter anchor, the frozen workflow executes the following
sequence:
\begin{enumerate}
  \item validate canonical daily observations and identify a persistent
  reported Method Code transition;
  \item construct a geographic candidate pool at distinct physical sites,
  retaining at most one POC per site by the pre-event ranking rule;
  \item record an explicit donor-insufficient disposition if fewer than three
  qualified physical donors exist;
  \item fit cross-site weights only on the 180-day pre-anchor calibration
  interval ending 15 days before the anchor;
  \item require at least two available donors and at least 30 retained
  observations in each calibration, pre-anchor, and post-anchor residual
  window;
  \item record an explicit input-window failure if those common conditions do
  not hold;
  \item for complete cases, compute fixed and nested diagnostic intervals,
  placebo diagnostics, donor sensitivity, and the rule-based evidence tier;
  \item preserve every row, including all failures and unavailable
  diagnostics, in the audit tables.
\end{enumerate}

This sequence is an implementation specification for the frozen audit. It
does not claim that any step confirms a physical mechanism.

\subsection{Case-study reconstruction contract}

The representative-case configuration is a display contract rather than a new
experiment. It fixes selection rules before rendering, binds public source
archives by their saved SHA-256 values, uses the frozen data-cleaning and
fixed-prior weight rules, and records the selected donor physical sites and
reconstruction tolerance. If a required archive is unavailable or mismatched,
the renderer fails rather than silently substituting a current public download.
No raw archive is distributed with this paper.

\subsection{Full perturbation-family matrix}

Table~\ref{tab:perturbation-metrics} contains every frozen method, metric, and
paired local/regional perturbation family. It is placed in the appendix for
readability; Figure~\ref{fig:perturbation-metrics} presents the same complete
matrix visually in the main text.

\begin{landscape}
{\small
\setlength{\tabcolsep}{4pt}
\renewcommand{\arraystretch}{1.06}
\begin{longtable}{@{}p{0.36\linewidth}rrrr@{}}
\caption{Perturbation-family-specific held-out synthetic metrics for all methods.}\label{tab:perturbation-metrics}\\
\toprule
Method & MAE & AUPRC & Macro-F1 & Regional FPR \\
\midrule
\endfirsthead
\multicolumn{5}{l}{\small\itshape Table \ref{tab:perturbation-metrics}, continued}\\
\toprule
Method & MAE & AUPRC & Macro-F1 & Regional FPR \\
\midrule
\endhead
\bottomrule
\multicolumn{5}{p{0.92\linewidth}}{\footnotesize\textit{Note.} Each family aggregates its target-only local perturbation and its matched target-and-donor regional variant. Consequently, regional FPR is assessed within the same paired family; the rows are not local-only effect estimates. N/A has the same meaning as in Table~\ref{tab:all-methods}.}\\
\endfoot
\addlinespace[2pt]
\multicolumn{5}{@{}l}{\textbf{Additive step}}\\
\addlinespace[1pt]
Standard SC & 0.07788 & 0.99946 & 0.93092 & 0.138 \\
MetaShift fixed & 0.08373 & 0.99695 & 0.87969 & 0.233 \\
MetaShift CV & 0.07680 & 0.99801 & 0.90161 & 0.193 \\
Nearest-neighbor DiD & 0.08680 & 0.99210 & 0.93614 & 0.105 \\
Bayesian mean shift & 0.22860 & 0.50497 & 0.41031 & 0.890 \\
Before-after median & 0.22774 & 0.50497 & 0.44778 & 0.808 \\
CUSUM & N/A & 0.50497 & 0.43785 & 0.833 \\
Rolling MAD & N/A & 0.50497 & 0.49862 & 0.553 \\
PELT & N/A & 0.50497 & 0.42673 & 0.858 \\
\addlinespace[2pt]
\multicolumn{5}{@{}l}{\textbf{Proportional step}}\\
\addlinespace[1pt]
Standard SC & 0.07635 & 0.85856 & 0.76319 & 0.188 \\
MetaShift fixed & 0.08016 & 0.84617 & 0.78293 & 0.090 \\
MetaShift CV & 0.07457 & 0.84296 & 0.78357 & 0.103 \\
Nearest-neighbor DiD & 0.08287 & 0.82011 & 0.70938 & 0.065 \\
Bayesian mean shift & 0.23444 & 0.50497 & 0.48849 & 0.350 \\
Before-after median & 0.22774 & 0.50497 & 0.46944 & 0.260 \\
CUSUM & N/A & 0.50497 & 0.48467 & 0.328 \\
Rolling MAD & N/A & 0.50497 & 0.40894 & 0.108 \\
PELT & N/A & 0.50497 & 0.48326 & 0.320 \\
\addlinespace[2pt]
\multicolumn{5}{@{}l}{\textbf{Gradual drift}}\\
\addlinespace[1pt]
Standard SC & 0.08321 & 0.99710 & 0.93985 & 0.110 \\
MetaShift fixed & 0.08551 & 0.99224 & 0.89419 & 0.193 \\
MetaShift CV & 0.07958 & 0.99521 & 0.92082 & 0.153 \\
Nearest-neighbor DiD & 0.08746 & 0.98570 & 0.93248 & 0.085 \\
Bayesian mean shift & 0.24325 & 0.50497 & 0.45055 & 0.800 \\
Before-after median & 0.23075 & 0.50497 & 0.47143 & 0.733 \\
CUSUM & N/A & 0.50497 & 0.45322 & 0.793 \\
Rolling MAD & N/A & 0.50497 & 0.49969 & 0.475 \\
PELT & N/A & 0.50497 & 0.49992 & 0.513 \\
\addlinespace[2pt]
\multicolumn{5}{@{}l}{\textbf{Temporary step}}\\
\addlinespace[1pt]
Standard SC & 0.16187 & 0.98113 & 0.91246 & 0.108 \\
MetaShift fixed & 0.16436 & 0.96551 & 0.88871 & 0.130 \\
MetaShift CV & 0.15875 & 0.96994 & 0.90372 & 0.113 \\
Nearest-neighbor DiD & 0.18885 & 0.93733 & 0.84114 & 0.065 \\
Bayesian mean shift & 0.28058 & 0.50497 & 0.49930 & 0.463 \\
Before-after median & 0.32696 & 0.50497 & 0.49614 & 0.588 \\
CUSUM & N/A & 0.50497 & 0.48686 & 0.660 \\
Rolling MAD & N/A & 0.50497 & 0.45747 & 0.780 \\
PELT & N/A & 0.50497 & 0.47750 & 0.708 \\
\addlinespace[2pt]
\multicolumn{5}{@{}l}{\textbf{Variance increase}}\\
\addlinespace[1pt]
Standard SC & N/A & 0.58680 & 0.48825 & 0.158 \\
MetaShift fixed & N/A & 0.57484 & 0.45372 & 0.105 \\
MetaShift CV & N/A & 0.57314 & 0.46902 & 0.113 \\
Nearest-neighbor DiD & N/A & 0.56157 & 0.42775 & 0.055 \\
Bayesian mean shift & N/A & 0.50056 & 0.45595 & 0.205 \\
Before-after median & N/A & 0.50741 & 0.42520 & 0.135 \\
CUSUM & N/A & 0.50470 & 0.48012 & 0.290 \\
Rolling MAD & N/A & 0.50191 & 0.39371 & 0.093 \\
PELT & N/A & 0.49172 & 0.45495 & 0.213 \\
\end{longtable}
}

\end{landscape}
